
\documentclass[11pt,a4paper]{article}
\usepackage[utf8]{inputenc}
\usepackage[T1]{fontenc}
\usepackage[french]{babel}
\usepackage{lmodern}
\usepackage{microtype}
\usepackage[a4paper,margin=1.65cm,headheight=15pt]{geometry}
\usepackage{amsmath,amssymb,mathtools}
\usepackage{array,tabularx,multicol,booktabs}
\usepackage{enumitem}
\usepackage[most]{tcolorbox}
\usepackage{xcolor}
\usepackage{tikz}
\usetikzlibrary{arrows.meta,calc,positioning}
\usepackage{pdflscape}
\usepackage{fancyhdr}
\usepackage{hyperref}
\usepackage{needspace}
\usepackage{pagecolor}

\definecolor{fondpage}{RGB}{247,248,250}
\definecolor{encre}{RGB}{31,35,40}
\definecolor{bleu}{RGB}{40,91,135}
\definecolor{vert}{RGB}{55,125,92}
\definecolor{orange}{RGB}{178,105,42}
\definecolor{violet}{RGB}{101,77,145}
\definecolor{rouge}{RGB}{170,72,72}
\definecolor{gris}{RGB}{86,91,99}
\definecolor{bleufonce}{RGB}{28,55,120}
\definecolor{bleuclair}{RGB}{238,243,255}
\definecolor{grisclair}{RGB}{245,245,245}
\definecolor{tableline}{RGB}{45,45,45}
\definecolor{softgray}{RGB}{246,247,249}
\definecolor{deepblue}{RGB}{25,62,115}
\definecolor{accent}{RGB}{20,82,130}

\hypersetup{colorlinks=true,linkcolor=bleufonce,urlcolor=bleufonce,
pdftitle={Parcours de revision -- Suites -- Premiere specialite},pdfauthor={}}

\setlength{\parindent}{0pt}
\setlength{\parskip}{0.25em}
\renewcommand{\arraystretch}{1.13}
\setlength{\tabcolsep}{4pt}
\setlist[itemize]{leftmargin=1.25em,itemsep=0.15em,topsep=0.2em}
\setlist[enumerate,1]{label=\textbf{\arabic*.}, leftmargin=1.05cm, itemsep=0.35em, topsep=0.2em}
\setlist[enumerate,2]{label=\textbf{\alph*.}, leftmargin=0.85cm, itemsep=0.25em, topsep=0.15em}

\pagestyle{fancy}
\fancyhf{}
\cfoot{\thepage}

\newcommand{\R}{\mathbb{R}}
\newcommand{\N}{\mathbb{N}}
\newcommand{\sourceq}[1]{ {\scriptsize\textcolor{bleufonce}{(site Première q#1)}}}
\newcommand{\code}[1]{\texttt{#1}}

\newcounter{question}
\newcounter{reponse}

\newenvironment{blocquestions}[1]{%
  \begin{tcolorbox}[enhanced,breakable,colback=bleuclair,colframe=bleufonce,boxrule=0.6pt,arc=2mm,left=2mm,right=2mm,top=2mm,bottom=1.5mm,title={\bfseries #1},coltitle=white,colbacktitle=bleufonce,before skip=3mm,after skip=3mm, fontupper=\large]
  \large
  \begin{enumerate}[leftmargin=*,label=\textbf{\arabic*.},itemsep=0.82em,topsep=0.5em]
  \setcounter{enumi}{\value{question}}
}{%
  \setcounter{question}{\value{enumi}}
  \end{enumerate}
  \end{tcolorbox}
}

\newenvironment{blocreponses}[1]{%
  \begin{tcolorbox}[enhanced,breakable,colback=bleuclair,colframe=bleufonce,boxrule=0.6pt,arc=2mm,left=2mm,right=2mm,top=2mm,bottom=1.5mm,title={\bfseries #1},coltitle=white,colbacktitle=bleufonce,before skip=3mm,after skip=3mm, fontupper=\large]
  \large
  \begin{enumerate}[leftmargin=*,label=\textbf{\arabic*.},itemsep=0.42em,topsep=0.5em]
  \setcounter{enumi}{\value{reponse}}
}{%
  \setcounter{reponse}{\value{enumi}}
  \end{enumerate}
  \end{tcolorbox}
}

\newtcolorbox{correctionbox}[1]{enhanced,breakable,colback=grisclair,colframe=black!55,boxrule=0.55pt,arc=2mm,left=2mm,right=2mm,top=2mm,bottom=1.5mm,title=\textbf{#1},coltitle=black,colbacktitle=black!12,before skip=2mm,after skip=2mm}

\newtcolorbox{methodebox}[1]{enhanced,breakable,colback=white,colframe=bleu!70!black,boxrule=0.55pt,arc=2mm,left=2mm,right=2mm,top=2mm,bottom=1.5mm,title=\textbf{#1},coltitle=white,colbacktitle=bleu!70!black,before skip=2mm,after skip=2mm}

\newcommand{\titreSujet}[2]{%
  \subsection*{#1}
  \addcontentsline{toc}{subsection}{#1}
  \textit{Référence / inspiration : #2}\par\medskip\hrule\medskip
}

\newcommand{\qcm}[4]{%
\begin{center}
\begin{tabularx}{0.96\linewidth}{@{}XX@{}}
\textbf{a.} #1 & \textbf{b.} #2\\[1mm]
\textbf{c.} #3 & \textbf{d.} #4
\end{tabularx}
\end{center}}

\tcbset{enhanced,arc=2mm,outer arc=2mm,boxrule=0.42pt,left=1.15mm,right=1.15mm,top=0.75mm,bottom=0.75mm,boxsep=0.35mm,before skip=0.8mm,after skip=0.8mm,fonttitle=\bfseries\footnotesize,coltitle=encre,before upper=\raggedright\fontsize{7.65}{8.15}\selectfont}
\newtcolorbox{fichebox}[3][]{colback=white,colframe=#2!72!black,colbacktitle=#2!18,coltitle=#2!50!black,borderline west={0.9mm}{0pt}{#2!78!black},title={#3},drop fuzzy shadow=black!5,#1}
\newcommand{\tagcol}[2]{\begin{tcolorbox}[enhanced,colback=#1!11,colframe=#1!45!black,boxrule=0.42pt,arc=2mm,left=1mm,right=1mm,top=0.45mm,bottom=0.45mm,before skip=0mm,after skip=0.85mm,fontupper=\bfseries\scriptsize,colupper=#1!55!black]\centering #2\end{tcolorbox}}
\newtcolorbox{panel}[1]{colback=fondpage,colframe=fondpage,boxrule=0pt,arc=0mm,left=0mm,right=0mm,top=0mm,bottom=0mm,height=168mm,valign=top,before skip=0mm,after skip=0mm}
\newcommand{\titrepage}{\begin{tcolorbox}[colback=white,colframe=bleu!70!black,boxrule=0.65pt,arc=3mm,left=2.5mm,right=2.5mm,top=1.15mm,bottom=1.15mm,before skip=0mm,after skip=1.4mm,drop fuzzy shadow=black!10]
{\Large\bfseries Parcours de révision -- Suites}\hfill {\normalsize Première spécialité -- sans calculatrice}
\end{tcolorbox}}


\newcommand{\miniAlgoNiemeTerme}{%
\begin{tcolorbox}[colback=white,colframe=gris!60!black,boxrule=0.35pt,arc=1mm,left=1.2mm,right=1mm,top=0.8mm,bottom=0.8mm,before skip=0.4mm,after skip=0.4mm]
\scriptsize
\begin{tabular}{@{}l@{}}
\code{u = u0}\\
\code{for k in range(n):}\\
\quad\code{u = f(u)}\\
\code{return u}
\end{tabular}
\end{tcolorbox}}

\newcommand{\miniAlgoSeuil}{%
\begin{tcolorbox}[colback=white,colframe=gris!60!black,boxrule=0.35pt,arc=1mm,left=1.2mm,right=1mm,top=0.8mm,bottom=0.8mm,before skip=0.4mm,after skip=0.4mm]
\scriptsize
\begin{tabular}{@{}l@{}}
\code{u = terme initial}\\
\code{n = 0}\\
\code{while condition:}\\
\quad\code{u = terme suivant}\\
\quad\code{n = n + 1}
\end{tabular}
\end{tcolorbox}}

\begin{document}
\begin{center}
{\LARGE\bfseries Parcours de révision -- Suites}\\[1mm]
{\large Première générale -- spécialité mathématiques}\\[1mm]
{\normalsize Sans calculatrice}
\end{center}
\vspace{2mm}
\tableofcontents
\clearpage

\phantomsection
\addcontentsline{toc}{section}{Partie 1 -- Récapitulatif de cours}
\newgeometry{margin=6.5mm}
\pagestyle{empty}
\begin{landscape}
\pagecolor{fondpage}\color{encre}
\thispagestyle{empty}

\fontsize{7.85}{8.45}\selectfont

\titrepage

\noindent
\begin{minipage}[t]{0.242\linewidth}
\tagcol{bleu}{1. DÉFINIR UNE SUITE}
\begin{panel}{bleu}

\begin{fichebox}{bleu}{Vocabulaire}
Une suite numérique est une fonction définie sur $\N$ ou sur une partie de $\N$.

Le terme de rang $n$ se note $u_n$.

Attention au premier rang : selon l'énoncé, la suite peut commencer à $u_0$ ou à $u_1$.
\end{fichebox}

\begin{fichebox}{vert}{Deux modes de définition}
\textbf{Forme explicite} : on donne $u_n$ en fonction de $n$.
\[
u_n=f(n).
\]
\textbf{Récurrence} : on donne un premier terme et une relation.
\[
u_0\text{ donné},\qquad u_{n+1}=f(u_n).
\]
\end{fichebox}

\begin{fichebox}{orange}{Calculer des termes}
Si $u_{n+1}$ dépend de $u_n$, on calcule les termes dans l'ordre.
\[
 u_0\longrightarrow u_1\longrightarrow u_2\longrightarrow \cdots
\]
Pour une formule explicite, on remplace directement $n$.
\end{fichebox}

\begin{fichebox}{rouge}{Algorithme : calculer $u_n$}
Si une suite commence à $u_0$ et vérifie $u_{n+1}=f(u_n)$, on répète $n$ fois le calcul du terme suivant.
\miniAlgoNiemeTerme
À la fin, la variable \code{u} contient le terme cherché.
\end{fichebox}

\end{panel}
\end{minipage}\hfill
\begin{minipage}[t]{0.242\linewidth}
\tagcol{vert}{2. SUITES ARITHMÉTIQUES}
\begin{panel}{vert}

\begin{fichebox}{vert}{Définition}
Une suite est arithmétique lorsqu'on ajoute toujours le même nombre.
\[
 u_{n+1}=u_n+r.
\]
Le réel $r$ est la \textbf{raison}.
\end{fichebox}

\begin{fichebox}{bleu}{Formule explicite}
Si le premier terme est $u_0$, alors :
\[
 u_n=u_0+nr.
\]
Si le premier terme est $u_p$, alors :
\[
 u_n=u_p+(n-p)r.
\]
\end{fichebox}

\begin{fichebox}{orange}{Variations}
Pour une suite arithmétique :
\[
\begin{array}{c|c}
r>0 & \text{croissante}\\
r=0 & \text{constante}\\
r<0 & \text{décroissante}
\end{array}
\]
\end{fichebox}

\begin{fichebox}{violet}{Somme}
Somme de termes consécutifs :
\[
\text{nombre de termes}\times\frac{\text{premier}+\text{dernier}}{2}.
\]
En particulier :
\[
1+2+\cdots+n=\frac{n(n+1)}2.
\]
\end{fichebox}

\end{panel}
\end{minipage}\hfill
\begin{minipage}[t]{0.242\linewidth}
\tagcol{orange}{3. SUITES GÉOMÉTRIQUES}
\begin{panel}{orange}

\begin{fichebox}{orange}{Définition}
Une suite est géométrique lorsqu'on multiplie toujours par le même nombre.
\[
 u_{n+1}=q u_n.
\]
Le réel $q$ est la \textbf{raison}.
\end{fichebox}

\begin{fichebox}{bleu}{Formule explicite}
Si le premier terme est $u_0$, alors :
\[
 u_n=u_0q^n.
\]
Si le premier terme est $u_p$, alors :
\[
 u_n=u_pq^{n-p}.
\]
\end{fichebox}

\begin{fichebox}{vert}{Variations}
Pour une suite géométrique à termes strictement positifs :
\[
\begin{array}{c|c}
q>1 & \text{croissante}\\
q=1 & \text{constante}\\
0<q<1 & \text{décroissante}
\end{array}
\]
\end{fichebox}

\begin{fichebox}{violet}{Somme}
Si $q\neq1$ :
\[
 u_0+u_1+\cdots+u_n
 =u_0\frac{1-q^{n+1}}{1-q}.
\]
Cas utile :
\[
1+q+\cdots+q^n=\frac{1-q^{n+1}}{1-q}.
\]
\end{fichebox}

\end{panel}
\end{minipage}\hfill
\begin{minipage}[t]{0.242\linewidth}
\tagcol{violet}{4. MÉTHODES ET MODÈLES}
\begin{panel}{violet}

\begin{fichebox}{violet}{Étudier les variations}
Méthode classique :
\[
 u_{n+1}-u_n.
\]
Si $u_{n+1}-u_n\geq0$ pour tout $n$, la suite est croissante.

Si $u_{n+1}-u_n\leq0$ pour tout $n$, elle est décroissante.
\end{fichebox}

\begin{fichebox}{bleu}{Taux d'évolution}
Une hausse de $t\%$ correspond au coefficient :
\[
1+\frac{t}{100}.
\]
Une baisse de $t\%$ correspond au coefficient :
\[
1-\frac{t}{100}.
\]
Taux constant $\Rightarrow$ suite géométrique.
\end{fichebox}

\begin{fichebox}{orange}{Algorithme de seuil}
Pour chercher le premier rang où un seuil est atteint, on utilise souvent une boucle \code{while}.
\miniAlgoSeuil
La variable $n$ compte le nombre de passages dans la boucle.
\end{fichebox}

\begin{fichebox}{rouge}{Choisir le bon modèle}
Accroissement constant $\Rightarrow$ suite arithmétique.

Taux constant $\Rightarrow$ suite géométrique.

Relation $u_{n+1}=a u_n+b$ avec $b\neq0$ : ce n'est en général ni arithmétique ni géométrique.
\end{fichebox}

\end{panel}
\end{minipage}

\end{landscape}
\nopagecolor
\restoregeometry
\pagestyle{fancy}
\clearpage
\section*{Partie 2 -- Questions flash}\addcontentsline{toc}{section}{Partie 2 -- Questions flash}
\begin{blocquestions}{Questions flash -- Suites}
  \item Qu'est-ce qu'une suite numérique ? \sourceq{1}
  \item Comment note-t-on en général le terme de rang $n$ d'une suite ? \sourceq{2}
  \item Quelle est la différence entre une définition explicite et une définition par récurrence ? \sourceq{3}
  \item Quand dit-on qu'une suite est arithmétique ? \sourceq{4}
  \item Comment s'appelle le réel $r$ dans une suite arithmétique ? \sourceq{5}
  \item Quelle est la formule explicite d'une suite arithmétique de premier terme $u_0$ et de raison $r$ ? \sourceq{6}
  \item Quand dit-on qu'une suite est géométrique ? \sourceq{7}
  \item Comment s'appelle le réel $q$ dans une suite géométrique ? \sourceq{8}
  \item Quelle est la formule explicite d'une suite géométrique de premier terme $u_0$ et de raison $q$ ? \sourceq{9}
  \item Comment étudie-t-on en général le sens de variation d'une suite ? \sourceq{10}
  \item Que montre-t-on pour prouver qu'une suite est croissante ? \sourceq{11}
  \item Que montre-t-on pour prouver qu'une suite est décroissante ? \sourceq{12}
  \item À quel nombre peut-on comparer $\dfrac{u_{n+1}}{u_n}$ pour étudier les variations ? \sourceq{13}
  \item Quel type de boucle Python utilise-t-on souvent pour chercher un seuil ? \sourceq{14}
  \item Quelle est la raison d'une suite arithmétique vérifiant $u_{n+1}=u_n-3$ ? \sourceq{15}
  \item Quelle est la raison d'une suite géométrique vérifiant $u_{n+1}=1{,}2u_n$ ? \sourceq{16}
  \item Que vaut $u_5$ si $u_n=2n+1$ ? \sourceq{17}
  \item Que vaut $u_3$ si $u_n=5\times2^n$ ? \sourceq{18}
  \item Si $u_0=7$ et $u_{n+1}=u_n+2$, que vaut $u_1$ ? \sourceq{19}
  \item Si $u_0=4$ et $u_{n+1}=3u_n$, que vaut $u_1$ ? \sourceq{20}
  \item Si $u_{n+1}-u_n\geq0$ pour tout $n$, quel est le sens de variation de $(u_n)$ ? \sourceq{21}
  \item Si tous les termes de $(u_n)$ sont positifs et si $\dfrac{u_{n+1}}{u_n}\leq1$, que peut-on dire ? \sourceq{22}
  \item Si une fonction Python contient \code{return u}, que renvoie-t-elle ? \sourceq{23}
  \item Dans un script de seuil, à quoi sert souvent \code{while} ? \sourceq{24}
  \item Qu'appelle-t-on une suite monotone ? \sourceq{25}
  \item Qu'appelle-t-on une suite constante ? \sourceq{26}
  \item Que vaut $u_2$ si $u_0=2$ et si $(u_n)$ est géométrique de raison $3$ ? \sourceq{27}
  \item Que vaut $u_3$ si $u_0=-1$ et si $(u_n)$ est arithmétique de raison $4$ ? \sourceq{28}
  \item Si tous les termes sont positifs et si $\dfrac{u_{n+1}}{u_n}=0{,}7$, quel est le sens de variation de $(u_n)$ ? \sourceq{29}
  \item Que vaut $u_4$ si $u_n=10-0{,}5n$ ? \sourceq{30}
  \item Si $(u_n)$ est géométrique, de premier terme positif et de raison $q>1$, quel est son sens de variation ? \sourceq{31}
  \item Si $(u_n)$ est géométrique, de premier terme positif et de raison $0<q<1$, quel est son sens de variation ? \sourceq{32}
  \item Donner la formule de la somme des $n+1$ premiers termes d'une suite arithmétique de premier terme $u_0$ et de dernier terme $u_n$. \sourceq{33}
  \item Donner la formule de la somme des $n+1$ premiers termes d'une suite géométrique de premier terme $u_0$ et de raison $q\neq1$. \sourceq{34}
  \item Quel type de suite modélise une évolution à accroissement constant ? \sourceq{35}
  \item Quel type de suite modélise une évolution à taux constant ? \sourceq{36}
\end{blocquestions}
\clearpage
\section*{Partie 3 -- QCM d'automatismes}\addcontentsline{toc}{section}{Partie 3 -- QCM d'automatismes}
Pour chaque question, une seule réponse est exacte. Aucune calculatrice n'est nécessaire.
\begin{enumerate}
\item Une suite arithmétique de raison $5$ vérifie :
\qcm{$u_{n+1}=5u_n$}{$u_{n+1}=u_n+5$}{$u_n=5^n$}{$u_{n+1}=u_n-5$}
\item Si $u_0=3$ et $u_{n+1}=2u_n$, alors $u_3$ vaut :
\qcm{6}{12}{24}{48}
\item La suite $u_n=7-2n$ est :
\qcm{arithmétique de raison $-2$}{géométrique de raison $-2$}{arithmétique de raison $7$}{constante}
\item Une baisse de $15\%$ correspond au coefficient multiplicateur :
\qcm{$1{,}15$}{$0{,}15$}{$0{,}85$}{$85$}
\item Si $(u_n)$ est géométrique, $u_0=5$ et $q=3$, alors :
\qcm{$u_n=5+3n$}{$u_n=3\times5^n$}{$u_n=5\times3^n$}{$u_n=15n$}
\item Pour chercher le premier rang à partir duquel $u_n\geq100$, on utilise souvent :
\qcm{une boucle \code{while}}{une commande \code{print} seulement}{une dérivée}{un tableau de signes}
\item Si $u_{n+1}-u_n<0$ pour tout $n$, alors la suite est :
\qcm{croissante}{décroissante}{géométrique}{positive}
\item La somme $1+2+\cdots+20$ vaut :
\qcm{$20^2$}{$\dfrac{20\times21}{2}$}{$2^{20}$}{$20+21$}
\end{enumerate}
\clearpage

\section*{Partie 4 -- Sujets type bac / E3C}
\addcontentsline{toc}{section}{Partie 4 -- Sujets type bac / E3C}

\titreSujet{Sujet 1 -- Grains de riz et somme géométrique}{inspiré du sujet zéro / specimen Première spécialité -- sans calculatrice}
On dispose d'une bande de $10$ cases. On place $3$ grains de riz sur la première case, puis on double le nombre de grains à chaque nouvelle case.

Pour tout entier $n$ compris entre $1$ et $10$, on note $u_n$ le nombre de grains de riz placés sur la case numéro $n$.
\begin{enumerate}
  \item Donner les valeurs de $u_1$, $u_2$ et $u_3$.
  \item Exprimer $u_{n+1}$ en fonction de $u_n$, pour $1\leq n\leq9$.
  \item Quelle est la nature de la suite $(u_n)$ ? Préciser sa raison.
  \item Donner l'expression de $u_n$ en fonction de $n$.
  \item Calculer le nombre de grains placés sur la dixième case.
  \item Calculer le nombre total de grains placés sur les $10$ cases.
  \item Expliquer pourquoi ce modèle correspond à une croissance géométrique et non à une croissance arithmétique.
\end{enumerate}

\titreSujet{Sujet 2 -- Prix d'un manteau en liquidation}{inspiré de sujets E3C Première sur pourcentages répétés -- sans calculatrice}
Un manteau coûte initialement $200$ euros. Chaque semaine, son prix diminue de $10\%$.

On note $u_n$ le prix du manteau, en euros, $n$ semaines après le début de la liquidation. Ainsi, $u_0=200$.
\begin{enumerate}
  \item Calculer $u_1$ et $u_2$.
  \item Justifier que la suite $(u_n)$ est géométrique et préciser sa raison.
  \item Donner l'expression de $u_n$ en fonction de $n$.
  \item Calculer le prix du manteau au bout de quatre semaines.
  \item Justifier le sens de variation de la suite $(u_n)$.
  \item Expliquer pourquoi il est faux de dire : « quatre semaines de baisse de $10\%$, c'est une baisse totale de $40\%$ ».
  \item Écrire, sans le calculer à la calculatrice, le montant total obtenu en additionnant les cinq premiers prix proposés : $u_0+u_1+u_2+u_3+u_4$.
\end{enumerate}

\titreSujet{Sujet 3 -- Suite récurrente et seuil algorithmique}{inspiré de sujets E3C Première sur suites et Python -- sans calculatrice}
On considère la suite $(u_n)$ définie par :
\[
 u_0=15,
 \qquad
 u_{n+1}=0{,}8u_n+1
\]
pour tout entier naturel $n$.
\begin{enumerate}
  \item Calculer $u_1$ et $u_2$.
  \item Montrer que si $u_n>5$, alors $u_{n+1}>5$.
  \item Montrer que, tant que $u_n>5$, on a $u_{n+1}<u_n$.
  \item Compléter le tableau suivant à la main. On ne demande pas de prolonger les calculs au-delà de $n=4$ :
\[
\begin{array}{c|ccccc}
n & 0 & 1 & 2 & 3 & 4\\ \hline
u_n & 15 & & & &
\end{array}
\]
  \item On considère l'algorithme Python suivant :
\[
\begin{array}{l}
\code{def seuil():}\\
\quad\code{u = 15}\\
\quad\code{n = 0}\\
\quad\code{while u > 6:}\\
\qquad\code{u = 0.8*u + 1}\\
\qquad\code{n = n + 1}\\
\quad\code{return n}
\end{array}
\]
On donne le tableau suivant, obtenu en poursuivant les calculs :
\[
\begin{array}{c|ccc}
n&9&10&11\\ \hline
u_n&6{,}34217728&6{,}073741824&5{,}8589934592
\end{array}
\]
À l'aide de ce tableau, déterminer ce que renvoie l'algorithme et interpréter le résultat dans le contexte de la suite.
  \item Modifier uniquement la condition de la boucle pour obtenir le premier rang $n$ tel que $u_n\leq7$.
  \item Expliquer pourquoi la suite $(u_n)$ n'est pas arithmétique et n'est pas géométrique.
\end{enumerate}

\end{document}
